\documentclass[11pt]{article}

\usepackage[margin=25mm]{geometry}

\usepackage{amsmath,amssymb,fontspec,unicode-math}

\IfFontExistsTF{DejaVu Serif}{\setmainfont{DejaVu Serif}}{\setmainfont{Latin Modern Roman}}

\IfFontExistsTF{Latin Modern Math}{\setmathfont{Latin Modern Math}}{\setmathfont{DejaVu Math TeX Gyre}}

\usepackage{xurl,hyperref}

\setlength{\parindent}{0pt}\setlength{\parskip}{6pt}

\begin{document}

\title{Dual-Flow: Transferable Multi-Target, Instance-Agnostic Attacks via In-the-wild Cascading Flow Optimization}

\author{ELVIA preliminary AI analysis}\date{21 October 2026}\maketitle

Preliminary assessment: Major.\par

Independent scientific review remains pending. Author adoption has not been requested.

Small ODE time steps do not control a fixed parameter step; an explicit smooth counterexample refutes Theorem 2's descent implication. The Morse cutoff proof has a separate local strict-growth gap. The ideal-flow statements and empirical attack results remain distinct.

\section*{Review scope}

Latest-version, claim-focused review of the cited author-paper passages and applicable current author-release paths. The stated counterexamples and scope controls were checked analytically; no author program, old verifier, training run or benchmark was executed.

\section*{Problem 1: Small time steps do not guarantee loss descent}

Source: Theorem 2, p.4; Algorithm 1, p.5; Proposition 4/Appendix A.2, pp.24–25.

Assessment: supported; impact: major

\subsection*{Source claim (paraphrase)}

Theorem 2 claims an update improves next-time extrapolated loss when the time decrement is sufficiently small and functions are smooth.

\subsection*{Analysis}

Choose constant velocity parameters \(\theta=\phi=0.02\), time \(t=0.5\), input \(0.49\), binary logits \((-10000(u-0.5)^2,0)\), and parameter step \(0.001\). The gradient is \(73.105857863\), so \(\theta^+=-0.053105857863\). Old next-time extrapolation remains \(0.49\); the updated extrapolation is \(0.5265529289315-0.073105857863\delta\). For every \(0<\delta\le0.1\), this moves farther from \(0.5\), increasing cross-entropy with clipping inactive. At \(\delta=0.05\), losses are \(1.3132616875\) and \(5.2482877367\). Arbitrarily small time decrements cannot repair an uncontrolled finite parameter step.

\subsection*{Finding and scope}

Theorem2 lacks a parameter-step condition needed for its claimed next-time descent; arbitrarily small time decrements do not repair this smooth finite-step counterexample.

\subsection*{Proposed correction}

Control the parameter learning rate and gradient alignment through a descent argument, as well as the ODE time discretization.

Proposed correction; independent verification and author adoption remain unestablished.

\section*{Problem 2: The cutoff construction does not give almost-everywhere strict growth}

Source: Proposition 1, p.4; Appendix A.1, pp.23–24.

Assessment: supported; impact: minor

\subsection*{Source claim (paraphrase)}

Proposition 1 requires a nondecreasing Morse-energy flow with strict growth almost everywhere.

\subsection*{Analysis}

The proof sets an extra cutoff to zero throughout neighborhoods of critical points, not only at the points. Every noncritical initial point in such an open frozen region has constant energy, contradicting almost-everywhere strictness. Nondecrease and existence of some suitable flow are not refuted.

\subsection*{Finding and scope}

The displayed cutoff construction does not establish almost-everywhere strict energy growth. Nondecrease and existence of a different suitable flow are not refuted.

\subsection*{Proposed correction}

An existing narrower construction uses \(X=\mu^2\nabla j\), with an interior-positive boundary-defining \(\mu\), avoiding the extra frozen critical neighborhoods. Its scope is a smooth bounded domain with finite-time boundary noncontact, not arbitrary hard-clipped image domains.

Proposed correction; independent verification and author adoption remain unestablished.

\section*{Limits}

Independent scientific adjudication remains pending.\par

No author code, model training or full benchmark evaluation was executed in this revision.\par

Version-specific source defects do not establish which code produced the published experiments.\par

Corrected benchmark scores, empirical effect sizes and author adoption remain unverified unless explicitly stated.\par

No new author-code execution, model spending, training, author contact or complete benchmark replication.\par

Historical numerical results are retained evidence rather than current reruns.\par

Local counterexamples do not invalidate every empirical result; missing historical run attribution remains unresolved.\par

\section*{Source identity}

NeurIPS 2025 official proceedings PDF (current publisher bytes, checked 2026-10-21)

\url{https://proceedings.neurips.cc/paper_files/paper/2025/file/6542b594cc49412ba4c0dd7e26d8a7ea-Paper-Conference.pdf}

Reviewed PDF SHA-256: \nolinkurl{cccf3754aa2e58c55c9262ed54f0fe287c1a3a767d7fa3da78f9c5f55d3617dd}

This is the same research work as the predecessor, matched by title, author identity and official publisher channel. Its freshly retrieved exact PDF bytes and real retrieval time are recorded. Current target arXiv metadata and the relevant available author-version passages were checked separately; background works were excluded. Original report/source history is unchanged.

\end{document}
